\begin{gather}\label{Rdef}
R(a_{+},a_{-},{\bf c};v,\hat{v}),\qquad a_{+}, a_{-},v,\hat{v}\in\C,\qquad a_{+}/a_{-}\notin (-\infty,0],\qquad {\bf c}\in\C^4,
\end{gather}
can be obtained as a limit of Spiridonov's 9-variable hyperbolic hypergeometric function~\cite{spir}. Their novel viewpoint leads to a third representation for the $R$-function. (See Proposition~4.20 and Theorem~4.21 in~\cite{bust} for the latter two representations.)

In this paper we are concerned with a 5-variable specialization of the $R$-function, def\/ined by
\begin{gather}\label{cRdef}
\cR(a_{+},a_{-},b;x,y)\equiv R(a_{+},a_{-},(b,0,0,0);x,y/2).
\end{gather}
Suitable discretizations of this function give rise to the $q$-Gegenbauer polynomials, whereas discretizations of the $R$-function yield the Askey--Wilson polynomials, cf.~I; moreover, the nonrelativistic limit of the $\cR$-function yields the conical function specialization of $_2F_1$. Hence it may be viewed as corresponding to the Lie algebra $A_1$, whereas the $R$-function can be tied in with $BC_1$.

The key new result of this paper concerning $\cR$ consists of the integral representation
\begin{gather}\label{key}
\cR(b;x,y)=\sqrt{\frac{\alpha}{2\pi}}\frac{G(2ib-ia)}{G(ib-ia)^2}\int_{\R}dz \frac{G(z+ (x-y)/2-ib/2)G(z- (x-y)/2-ib/2)}{G(z+ (x+y)/2+ib/2)G(z- (x+y)/2+ib/2)}.
\!\!\!
\end{gather}
Here and throughout the paper we use parameters
\begin{gather}\label{aconv}
\alpha\equiv 2\pi/a_+a_-,\qquad a\equiv (a_++a_-)/2,
\end{gather}
$G(\aaa;z)$ is the hyperbolic gamma function (cf.~Appendix~\ref{appendixA}), and  the dependence on~$a_{+}$,~$a_{-}$ is suppressed. (We shall often do this when no confusion can arise.) Furthermore, in~\eqref{key} we choose at f\/irst
\begin{gather}\label{first}
(\aaa,b,x,y)\in(0,\infty)^2\times(0,2a)\times\R^2.
\end{gather}

By contrast to the previous three integral representations following from I, \cite{bult} and~\cite{bust}, the integrand in~\eqref{key}
involves only four hyperbolic gamma functions. We also obtain several closely related representations that involve in addition plane waves, cf.~\eqref{Ri}--\eqref{Rv}. As will transpire in Section~\ref{section3}, upon using the f\/irst one~\eqref{key} of these novel representations (which we dub `minimal' representations) to introduce the $\cR$-function, it is possible to rederive in a more transparent and self-contained way a great many features that also follow upon specialization of the $R$-function theory, developed not only in~I, II and~III, but also in our later papers~\cite{anne} and~\cite{quad}. Moreover, special cases and limits of the $\cR$-function are far more easily obtained from the minimal representations than from the original integral representation of I or from the alternative representations following from~\cite{bult} and~\cite{bust}. (The integrands of these earlier representations involve at least eight hyperbolic gamma function factors.)

A survey of the results of I--III and~\cite{anne} can be found in~\cite{Rsurv}, but the def\/inition~\eqref{cRdef} of the $A_1$-analog of the ($BC_1$) $R$-function dates back to the more recent paper~\cite{quad}. In Section~\ref{section2}, we review in particular the pertinent results from~\cite{quad}. However, we have occasion to  add a lot more information that follows by specializing previous f\/indings concerning the $R$-function and related functions to their $A_1$ counterparts. This includes the asymptotic behavior and Hilbert space properties obtained in~II and~III, resp., which are adapted to the $A_1$ setting in Subsection~\ref{section2.2}, and the parameter shifts obtained in~\cite{anne}, which we focus on in~Subsection~\ref{section2.4}.  Moreover, in~\eqref{cRM} we detail the connection of the renormalized function
\begin{gather}\label{cRr}
\cR_r(\aaa,b;x,y)\equiv  \frac{G(ib-ia)}{G(2ib-ia)}\cR(\aaa,b;x,y),
\end{gather}
to the $A_1$ type functions $M(ma_{+}+na_{-};x,y)$, $m,n\in\Z$, which featured in our previous papers~\cite{GLFII} and~\cite{Hilb}. We present the proof of~\eqref{cRM} in~Subsection~\ref{section2.3}, together with various corollaries.

Altogether, Section~\ref{section2} invokes a considerable amount of information from our previous work. We have attempted to sketch this in such a way that the reader need only consult the pertinent papers for quite technical aspects (in case of doubt and/or inclination, of course). Even so, it is probably advisable to skim through Section~\ref{section2} at f\/irst reading, referring back to it when the need arises.

By contrast, Section~\ref{section3} (combined with Appendices~A and~C) is largely self-contained. Its starting point is a hyperbolic functional identity that f\/irst arose as a specialization of elliptic functional identities expressing the relation of certain Hilbert--Schmidt integral kernels to the elliptic $BC_1$ relativistic Calogero--Moser dif\/ference operators introduced by van Diejen~\cite{diej}. We need not invoke these identities (which can be found in~\cite{HSI}, cf.~also~\cite{noum}), since the relevant hyperbolic version is quite easily proved directly. The key point is that the hyperbolic identities can be rewritten in terms of two pairs of hyperbolic $A_1$-type relativistic Calogero--Moser dif\/fe\-rence operators $A_{\pm}(b_j;x)$, $j=1,2$, with distinct couplings $b_1$, $b_2$. The dif\/ference operators are given by
\begin{gather}\label{Apm}
A_{\de}(b;x)\equiv \frac{s_{\de}(x-ib)}{s_{\de}(x)}T^x_{ia_{-\de}}+\frac{s_{\de}(x+ib)}{s_{\de}(x)}T^x_{-ia_{-\de}},\qquad \de=+,-.
\end{gather}
Here, the translations are def\/ined on analytic functions by
\begin{gather}\label{Tdef}
(T^z_cf)(z)\equiv f(z-c),\qquad c\in\C^{*}.
\end{gather}
Also, throughout this paper we use the abbreviations
\begin{gather}\label{denot}
s_{\de}(z)=\sinh(\pi z/a_{\de}),\! \qquad c_{\de}(z)=\cosh(\pi z/a_{\de}),\! \qquad e_{\de}(z)=\exp(\pi z/a_{\de}),\! \qquad \de=+,-.\!\!
\end{gather}
Choosing $b_1=b$, $b_2=0$, an auxiliary function $B(b;x,y)$ can be def\/ined that satisf\/ies the eigenvalue equations (at f\/irst under certain restrictions on the $B$-arguments)
\begin{gather}\label{AdeB}
A_{\de}(b;x)B(b;x,y)=2c_{\de}(y)B(b;x,y),\qquad \de=+,-.
\end{gather}
 More specif\/ically, the function $B(b;x,y)$ is the Fourier transform of the hyperbolic kernel function, which is a product of four hyperbolic gamma functions.

When we write the integrand of the integral def\/ining $B$ as a product of two factors that involve only two hyperbolic gamma functions, we can use the Plancherel relation and the explicit Fourier transform formula for factors of this type (derived in Appendix~\ref{appendixC}) to obtain two new integral representations for $B$.  In particular, this leads to a function $C(b;x,y)$ given by
\begin{gather}\label{defC}
C(b;x,y)\equiv \sqrt{\frac{\alpha}{2\pi}}\int_{\R}dz \frac{G(z+ (x-y)/2-ib/2)G(z- (x-y)/2-ib/2)}{G(z+ (x+y)/2+ib/2)G(z- (x+y)/2+ib/2)}.
\end{gather}

Comparing~\eqref{defC}, \eqref{key} and~\eqref{cRr}, we read of\/f
\begin{gather}\label{cRC}
\cR_r(b;x,y)=G(ib-ia)C(b;x,y).
\end{gather}
However, we need a further study of the $C$-function~\eqref{defC} to arrive at a proof of this relation to the function $\cR_r$, as def\/ined originally by~\eqref{cRr} and~\eqref{cRdef}. Indeed, as already alluded to below~\eqref{first}, we can use~\eqref{defC} as a starting point to derive many features that $C$ and $\cR_r$ have in common.

In particular, the general analysis in~Appendix~B of~I can be applied to the integral on the r.h.s.\ of~\eqref{defC}, which yields a complete elucidation of the behavior of $C(b;x,y)$ under meromorphic continuation.
Moreover, via the A$\De$Es (analytic dif\/ference equations)~\eqref{AdeB} and the manifest invariance of $C$ under interchanging $x$ and $y$, it follows that
$C(b;x,y)$ is a joint eigenfunction of the four A$\De$Os (analytic dif\/ference operators)
\begin{gather}\label{4A}
A_{+}(b;x),\qquad A_{-}(b;x),\qquad A_{+}(b;y),\qquad
A_{-}(b;y),
\end{gather}
with eigenvalues
\begin{gather}\label{4ev}
2c_{+}(y),\qquad 2c_{-}(y),\qquad
2c_{+}(x),\qquad 2c_{-}(x),
\end{gather}
resp. This is also the case for $\cR_r(b;x,y)$ and, moreover, the equality~\eqref{cRC} can be shown for the special case $y=ib$ by a further application of Appendix~\ref{appendixC}. The general case then follows by a uniqueness argument already used in Subsection~\ref{section2.3}.

We reconsider the special $b$-values
\begin{gather}\label{bmn}
b_{mn}\equiv ma_{+}+na_{-},\qquad m,n\in\Z,
\end{gather}
in Subsection~\ref{section4.1}, inasmuch as they satisfy $b_{mn}\in(0,2a)$.  Indeed, the new Fourier transform representations in Section~\ref{section3} are only well def\/ined for $b\in(0,2a)$, but they can be explicitly evaluated by a residue calculation when $b$ is of this form. The key point is that the $G$-ratios in the integrand can then be written in terms of the hyperbolic cosines $c_{\pm}(z)$ by using the $G$-A$\De$Es~\eqref{Gades}. In principle, this yields again the functions $M(b_{mn};x,y)$ from~\cite{GLFII}, but we have not tried to push through a direct equality proof (as opposed to appealing to uniqueness).
